\begin{table}[htbp]
\centering
\caption{Rendimiento Promedio por Categor\'ia de Modelo}
\label{tab:category_comparison}
\begin{tabular}{@{}lrrrrr@{}}
\toprule
\textbf{Categor\'ia} & \textbf{N} & \textbf{Ret. Prom.} & \textbf{Mejor Ret.} & \textbf{Sharpe Prom.} & \textbf{Mejor Modelo} \\
\midrule
ML cl\'asico & 4 & \textcolor{ForestGreen}{106.4\%} & 252.5\% & 0.403 & Ridge \\
Deep learning & 9 & \textcolor{ForestGreen}{72.6\%} & 395.4\% & 0.215 & LSTM\_Attention \\
Boosting & 4 & \textcolor{ForestGreen}{36.2\%} & 66.3\% & 0.217 & GradientBoosting \\
Series temporales & 6 & \textcolor{ForestGreen}{20.8\%} & 44.4\% & 0.134 & Theta \\
\bottomrule
\end{tabular}
\begin{tablenotes}
\small
\item Nota: N = n\'umero de modelos en la categor\'ia. ML cl\'asico = Ridge/Lasso/ElasticNet/RandomForest; Boosting = GradientBoosting (sklearn) y las librer\'ias XGBoost/LightGBM/CatBoost; Series temporales = AutoARIMA/ETS/Theta/SeasonalNaive/Prophet/GARCH; Deep learning = redes Darts (DLinear/N-BEATS/N-HiTS/TCN/TFT) y arquitecturas \textit{custom} en PyTorch (CNN-LSTM/LSTM+Attention/BiLSTM/BiGRU).
\end{tablenotes}
\end{table}
